\honest{Dark Side Epistemology}{Eliezer Yudkowsky}{2008}

``If you once tell a lie, the truth is ever after your enemy.''

I have said before that lies are contagious. Tell a geologist that a driveway pebble came from a beach, and the geologist may know better: a water-worn pebble does not look like frozen lava. A truth that seems arbitrary to one mind may be pinned down by a dozen links for a mind that knows more. To a creationist, design versus natural selection may sound like a choice of sports team. \nb{Yudkowsky's own footnote reports a geologist's reply that most driveway pebbles are taken from beaches, so this lie would pass. The footnote calls it a ``Case in point.''}

To a biologist, arguing that humans were designed would mean lying about the retina, the brain, and ``the proteins bound together by weak van der Waals forces instead of strong covalent bonds.'' \nb{The retina is a standard example in arguments from poor design. The post gives no reason why weak bonds would count against design; proteins usually change shape to do their work.} So most creationists lie about evolutionary theory instead, and then about the rules of science, such as what ``theory'' means. The lies spread from specific facts to general laws to the rules of reasoning. \nb{The post gives no evidence that most creationists know their claims to be false, and near the end it says that most who repeat bad epistemology are ``more duped than duplicitous.'' The logical point, that any belief can be kept if enough other beliefs are revised, even the laws of logic, was made by W. V. Quine in 1951.}

The same happens when you lie to yourself. I write an exchange with a believer in a dragon in the garage. Asked for evidence, the believer says dragons are exempt. Told why beliefs require evidence (a map of a city drawn at random in your living room will almost certainly be wrong, and a map of a dragon is no different), the believer says that ``almost certainly'' leaves a chance. Or the believer suppresses a dawning suspicion, or puts the dragon in ``a separate magisterium.'' A false belief can be corrected if you get over it when you find the mistake. What is dangerous is a belief you think must be protected as a belief, whether or not you actually believe it. ``A single Lie That Must Be Protected can block someone's progress into advanced rationality.'' \nb{The post names no one whose progress was blocked this way. Its support is the exchange the author wrote.}

Rational beliefs, like the world, are tied together by general laws more tightly than the untrained suspect. Think of all the connected truths you would have to refuse to know in order to deny evolution or heliocentrism. One self-deception can then block the whole higher level of truth-seeking, and put in its place rules of anti-thought and general justifications for believing what is false. Steven Kaas said that promoting inaccurate beliefs is sabotage, like slashing tires. Giving someone a belief they must protect is like giving them a lobotomy. To protect a lie you must deny that beliefs require evidence, that maps should reflect territories, and finally that truth is good. ``Thus comes into being the Dark Side.''

I worry that people are not wary enough of meeting systematically bad epistemology. Some popular advice on how to think was made by rationalists, and some ``invented to protect a lie or self-deception.'' My example: ``Everyone has a right to their own opinion.'' Was it said to protect a truth, or to protect from the truth? I do not look into where it came from. \nb{In its political sense it is a right to hold opinions ``without interference'' (Universal Declaration of Human Rights, Article 19), a claim about coercion, not about evidence. The question also judges the proverb by the purpose of whoever first said it, which sits uneasily with the rule the post states three paragraphs later.} But how else? Where lies are protected, bad epistemology will arise to fight the good.

Most who repeat this ``Deep Wisdom'' are ``more duped than duplicitous \ldots\ I think.'' And do not use my idea as an all-purpose rebuttal: ``you have to refute the proposition for itself, not by accusing its inventor of bad intentions.''

The Dark Side is out there, fear is the path to it, and there are Sith Lords as well as Jedi. I end by asking readers to list more ideas ``spawned by the Dark Side.''
